\chapter*{Abstract}
\addcontentsline{toc}{chapter}{Summary}

Capturing dynamic processes happening inside live cells is difficult due to the fact that
many subcellular structures are completely hidden beyond diffraction blur.
Super-resolution optical fluctuation imaging (SOFI) surpasses the optical diffraction
limit using only a conventional widefield microscope, but it requires hundreds of frames
and offline cumulant computation, making it too slow for real-time live-cell imaging.
Recent work (RESURF) showed that a recurrent deep network, SOFI-MISRGRU, can approximate a
second-order SOFI reconstruction from as few as eight frames at real-time speed. Its
ConvGRU fusion is, however, inherently sequential and aggregates frames through an
unnormalised summation that does not generalise beyond the training frame count.

This thesis investigates whether a Transformer that fuses all frames in parallel through
self-attention (TRMISR) can improve upon recurrent fusion in super-resolution quality 
across various signal-to-noise and blinking conditions, and at what inference-latency cost.
Secondly, we modify the SOFI-MISRGRU into a streaming architecture that 
emits one super-resolved image per incoming frame. 
Both models are adapted to the non-integer $2n{-}7$ SOFI upscaling and evaluated on a
synthetic microtubule dataset spanning twelve blinking and signal-to-noise conditions.
The models are using the resolution-scaled Pearson correlation and 
decorrelation analysis to match the RESURF evaluation.

TR-MISR matches the 20-frame MISRGRU baseline while holding its quality under low SNR,
slow-blinking acquisitions where MISRGRU degrades and hallucinates some filaments.
Because cross-attention normalises over the input tokens, an eight-frame-trained TRMISR can
be evaluated on longer sequences without retraining, a property the GRU lacks. Our modified 
streaming MISRGRU architecture, reaches the batch baseline by the twentieth accumulated frame 
and keeps improving. Model augmented with a deformable phase-space alignment module, and trained using motion augmentation
pipeline that simulates moving structures further
sustains reconstruction quality under inter-frame sample motion up to speeds of 5 $\mu$m per second.
By emitting one image per incoming frame instead of reprocessing the whole window, the
streaming model cuts the per-output latency from 170\,ms to 17\,ms, a $10\times$ speed-up over
the original RESURF work.
Together these results bring deep-learning super-resolution microscopy closer to real-time use.